% Generated draft: requires mathematical and source review.
\documentclass[11pt]{article}
\usepackage[margin=25mm]{geometry}
\usepackage{fontspec}
\setmainfont{Latin Modern Roman}
\usepackage{amsmath,amssymb,mathtools,unicode-math}
\setmathfont{Latin Modern Math}
\usepackage{enumitem,needspace,xurl,hyperref}
\setlength{\parindent}{0pt}
\setlength{\parskip}{5pt}
\setlength{\emergencystretch}{3em}
\newcommand{\N}{\mathbb{N}}
\newcommand{\Z}{\mathbb{Z}}
\newcommand{\Q}{\mathbb{Q}}
\newcommand{\R}{\mathbb{R}}
\newcommand{\C}{\mathbb{C}}
\newcommand{\F}{\mathbb{F}}
\newcommand{\A}{\mathbb{A}}
\newcommand{\PP}{\mathbb{P}}
\newcommand{\E}{\mathbb{E}}
\newcommand{\eps}{\varepsilon}
\newcommand{\dd}{\,\mathrm{d}}
\begin{document}
\section*{UnsolvedMath ID: 800946}
\textbf{Source problem code:} catalogue entry
\subsection{Definitions and mathematical statement}
Let \((M,g)\) be a closed surface of genus at least \(2\), equipped with a Riemannian metric whose Gaussian curvature \(K_g\) satisfies \(K_g\leq 0\) everywhere. Here closed means compact and without boundary. For \(L>0\), define
\[
Z_g(L)=\#\left\{\alpha\in [S^1,M]_{\mathrm{free}}:\begin{array}{l}
\text{there exists a nonconstant closed }g\text{-geodesic }\gamma\in\alpha,\\
\ell_g(\gamma)\leq L\text{ and }K_g(\gamma(t))=0\text{ for every }t\in S^1
\end{array}\right\},
\]
where \([S^1,M]_{\mathrm{free}}\) denotes free homotopy classes of maps from the circle to \(M\), and \(\ell_g(\gamma)\) is the length of the parametrized closed geodesic, including traversal multiplicity. Thus each free homotopy class is counted once, regardless of how many qualifying representatives it contains.

The question attributed to Knieper is whether, for every such \((M,g)\), there exists a finite constant \(C=C(M,g)\geq 0\) such that
\[
Z_g(L)\leq C L^2\qquad\text{for every }L>0.
\]
The source does not restrict the count to primitive geodesics. The displayed definition therefore includes classes represented by iterates and uses ordinary free homotopy of parametrized loops, without additionally identifying a class with its inverse. These counting conventions require confirmation against the intended source convention. The source also does not explicitly specify orientability or distinguish an eventual quadratic bound from a bound for every positive \(L\); the display states the latter interpretation.
\subsection{Short English statement}
On each closed surface of genus at least \(2\) with nonpositive curvature, count the free homotopy classes admitting a closed geodesic of length at most \(L\) along which Gaussian curvature vanishes everywhere. Knieper asks whether this count is bounded by \(C L^2\), with \(C\) allowed to depend on the surface and metric. The source leaves some counting and domain conventions implicit.
\subsection{Sources}
\begin{itemize}
\item Local UnsolvedMath registry, permanent ID \(800946\), source problem code \texttt{catalogue entry}, title: \emph{Open Problems and Questions about Geodesics --- Question 26}. The supplied record labels the problem open; this is a registry label, not an independently verified status assessment.
\item Burns and Matveev, \emph{Open Problems and Questions about Geodesics} (2019), Question 26, attributed in the supplied statement to Knieper. Supplied source URL: \url{https://arxiv.org/abs/1308.5417}. The registry lists Burns and Matveev as proposers and gives proposed year 2019.
\item Supplied provenance records access on 2026-10-16 and extraction by \texttt{arxiv\_latex}. This access information is supplied provenance; no independent source visit or literature verification is asserted here.
\end{itemize}
\end{document}
